\section{Method}
\label{sec:method}

\subsection{Forward dynamics}
\label{sec:forward-dynamics}

\paragraph{Oscillatory synapse and real-valued soma.}
For neuron $i$, let the complex synaptic state be
\begin{equation}
    z_i(t)=I_i(t)+\mathrm{i}Q_i(t),
\end{equation}
where only $I_i=\operatorname{Re}[z_i]$ drives the membrane voltage. A spike
from presynaptic neuron $j$ is delivered at its final spike time $t^s_{jk}$
and produces a quadrature jump of size $-w_{ij}$. The compact complex dynamics
are
\begin{equation}
    \dot z_i(t)=\left(-\alpha+\mathrm{i}\omega_i\right)z_i(t)
    -\mathrm{i}\sum_j w_{ij}\sum_k\delta\!\left(t-t^s_{jk}\right).
    \label{eq:tes-complex-synapse}
\end{equation}
Thus $1/\tau_I$ in an exponential EventProp synapse is replaced by the pair
$(\alpha,\omega_i)$. The present implementation fixes $\alpha$ and $\tau_V$
and assigns each neuron to one of two non-learned frequency groups with periods
$8\pi\,\mathrm{ms}$ and $12\pi\,\mathrm{ms}$, corresponding to
$\omega_{\mathrm{fast}}=1/4\,\mathrm{rad\,ms}^{-1}$ and
$\omega_{\mathrm{slow}}=1/6\,\mathrm{rad\,ms}^{-1}$, with
$\alpha=1/40\,\mathrm{ms}^{-1}$ and $\tau_V=20\,\mathrm{ms}$. The first
$\lceil N_h/2\rceil$ hidden neurons use the fast group, the remaining hidden
neurons use the slow group, and every readout uses the slow group. The synaptic
state decays with envelope
$\exp(-\alpha t)$ and rotates at angular frequency $\omega_i$.

And the soma follows the EventProp LIF convention
\begin{equation}
    \tau_V\dot V_i(t)=-V_i(t)+I_i(t).
    \label{eq:tes-voltage}
\end{equation}
Equivalently, between events,
\begin{equation}
    \frac{\mathrm d}{\mathrm dt}
    \begin{bmatrix}V_i\\ I_i\\ Q_i\end{bmatrix}
    =
    \underbrace{\begin{bmatrix}
        -1/\tau_V & 1/\tau_V & 0\\
        0 & -\alpha & -\omega_i\\
        0 & \omega_i & -\alpha
    \end{bmatrix}}_{A_i}
    \begin{bmatrix}V_i\\ I_i\\ Q_i\end{bmatrix}.
    \label{eq:tes-real-flow}
\end{equation}
At an arrival event,
\begin{equation}
    I_i\!\left((t^s_{jk})^+\right)
      =I_i\!\left((t^s_{jk})^-\right),
    \qquad
    Q_i\!\left((t^s_{jk})^+\right)
      =Q_i\!\left((t^s_{jk})^-\right)-w_{ij}.
    \label{eq:tes-arrival-jump}
\end{equation}
The impulse is applied to the quadrature state. Starting from zero state, one
arrival produces
$I_i(t)=w_{ij}e^{-\alpha(t-t^s_{jk})}
\sin(\omega_i(t-t^s_{jk}))$ for $t>t^s_{jk}$, so the current seen by the
soma is a damped, sign-changing response. Importantly, $I_i$ and therefore
$\dot V_i$ remain continuous at the arrival even though $Q_i$ jumps.

For a unit-weight arrival at time zero, writing $\beta=1/\tau_V$, the voltage
kernel is
\begin{equation}
 V(t)=\frac{e^{-\alpha t}
 \left[(\beta-\alpha)\sin(\omega t)-\omega\cos(\omega t)
 +\omega e^{-(\beta-\alpha)t}\right]}
 {\tau_V\left[(\beta-\alpha)^2+\omega^2\right]}.
 \label{eq:tes-unit-psp}
\end{equation}
Initialization matches the first positive-lobe peak to the $0.25$ unit peak of
the historical vanilla exponential-current reference. The fast group has
unscaled peak $0.2591420990$ at $11.1676331128\,\mathrm{ms}$ and gain
$0.9647216756$; the slow group has unscaled peak $0.3202711116$ at
$15.8712273368\,\mathrm{ms}$ and gain $0.7805886666$. These gains compensate
first-lobe excitability but do not remove oscillatory curvature. Moreover, the
$50\,\mathrm{ms}$ horizon exposes only finitely many later lobes, so temporal
expressiveness claims are limited to crossings that occur inside this window.

\paragraph{Threshold, reset, and state-dependent delay.}
Neuron $i$ reaches threshold at the $k$th upward crossing
\begin{equation}
    h_i\!\left(\mathbf{x}(t^\theta_{ik})\right)
    =V_i(t^\theta_{ik})-\vartheta_i=0,
    \qquad
    \dot V^\theta_{ik}>0,
    \label{eq:tes-threshold}
\end{equation}
where $\mathbf{x}$ denotes the state of the complete network and
$\dot V^\theta_{ik}=
\dot V_i((t^\theta_{ik})^-)$ is the crossing slope. The membrane is reset
immediately at $t^\theta_{ik}$. With configured reset drop $R_i>0$, we use the
subtractive reset
\begin{equation}
    V_i\!\left((t^\theta_{ik})^+\right)
      =V_i\!\left((t^\theta_{ik})^-\right)-R_i,
    \label{eq:tes-reset}
\end{equation}
with $I_i$ and $Q_i$ unchanged. A hard reset can be used by replacing this map
in the event Jacobian below.

Before applying the reset, $\dot V^\theta_{ik}$ is obtained directly by
evaluating Equation~\eqref{eq:tes-voltage} at the threshold event and no
numerical differentiation of a voltage trace is required. This slope determines
the axonal delay
\begin{equation}
    \mathfrak{d}_{ik}
      =d_{+}\sigma\!\left(c_1-c_0\dot V^\theta_{ik}\right).
    \label{eq:tes-delay-method}
\end{equation}
The spike arriving time is
\begin{equation}
    t^s_{ik}=t^\theta_{ik}+\mathfrak{d}_{ik}.
    \label{eq:tes-final-spike-time}
\end{equation}
Since $c_0>0$, a sharper crossing yields a shorter delay. The parameter $c_1$
sets the midpoint crossing slope $c_1/c_0$ and controls both the maximum delay
value $d_{+}\sigma(c_1)$ and the local derivative there.
If hardware requires a
strictly positive $d_{\min}$, the sigmoid output can be affinely rescaled to
$[d_{-},d_{+}]$.

For an isolated contribution delivered to postsynaptic neuron $m$, this
temporal shift corresponds to the phase offset
    $\Delta\phi_{m\leftarrow i,k}=\omega_m\mathfrak{d}_{ik}.
$
The delay therefore selects the phase at which the event enters the
postsynaptic oscillatory filter.

\paragraph{Event-driven forward pass.}
Between successive threshold and arrival events, Equation~\eqref{eq:tes-real-flow} is propagated exactly with a matrix exponential.  When a threshold event occurs, the simulator records the pre-reset slope, schedules the corresponding arrival at $t^s_{ik}$, and resets the source voltage. When a scheduled event occurs, Equation~\eqref{eq:tes-arrival-jump} is applied to every postsynaptic target. 

\paragraph{Layerwise feed-forward continuous evaluation.}
For the input--hidden--readout topology used in the present Yin--Yang
checkpoint, neurons within a layer are conditionally independent.  We
therefore solve every upward threshold crossing of each hidden neuron
independently from the exact external-arrival sequence.  After every crossing,
the configured subtractive reset is applied, an unquantized slope-dependent
arrival is scheduled, and the same hidden stream continues to the trial
horizon.  The merged hidden-arrival sequence then drives each readout stream,
which is likewise allowed to cross and reset repeatedly in full-horizon mode.
This layerwise decomposition does not require a global earliest-threshold
search within the hidden layer.  A hidden neuron without a physical threshold
crossing remains silent and emits no event.  Silent readouts are also absent
from the physical event tape, while only the first physical readout threshold
enters the squared-hinge objective.
For the active objective, let $P$ denote
the physical readouts.  If $P$ is nonempty, every silent readout is assigned
the report-only effective time

\begin{equation}
    \widetilde t_f=\max_{r\in P}t_r+2\,\mathrm{ms}.
    \label{eq:tes-fake-readout-time}
\end{equation}

If all readouts are silent, their report-only times remain $T$.  These values
are explicitly marked as censored: they are not threshold crossings, they do
not schedule arrivals, and they create no backward event.

The full continuous mode retains all physical events through the requested
horizon.  The optional early-stop mode instead terminates after every readout
has a first physical spike, or after the remaining readouts are certified
silent.  The latter certificate is applied only after all upstream arrival
streams are closed and uses the strict sufficient condition
$\max\{V,\sqrt{I^2+Q^2}\}<\vartheta$.  Full and early modes therefore have the
same effective first readouts and gradients, although their retained event
tapes and work counters can differ.  The grid-resolved solver uses the same
model and event signs, but places crossings at grid endpoints and quantizes a
hidden delay upward to an integer number of time steps.

\subsection{TeS backward equations}
\label{sec:tes-backward}

We consider an EventProp loss consisting of a spike-time term and a
voltage-dependent running term,
\begin{equation}
    \mathcal{L}
      =l_{\mathrm p}\!\left(\{t^s_{ik}\}\right)
       +\int_0^T l_V(\mathbf V(t),t)\,\mathrm dt.
    \label{eq:tes-loss}
\end{equation}
The adjoints are propagated from $T$ to $0$, and a prime denotes the reverse
time derivative, $'=-\mathrm d/\mathrm dt$. Between events they obey
\begin{align}
    \tau_V\lambda_{V_i}'
      &=-\lambda_{V_i}-\frac{\partial l_V}{\partial V_i},
      \label{eq:tes-adjoint-v}\\
    \lambda_{I_i}'
      &=\lambda_{V_i}-\alpha\lambda_{I_i}
        +\omega_i\lambda_{Q_i},
      \label{eq:tes-adjoint-i}\\
    \lambda_{Q_i}'
      &=-\omega_i\lambda_{I_i}-\alpha\lambda_{Q_i}.
      \label{eq:tes-adjoint-q}
\end{align}
For Equation~\eqref{eq:tes-loss}, the terminal conditions are zero. The pair
$(\lambda_{I_i},\lambda_{Q_i})$ therefore transports error backward through
the transpose oscillatory dynamics, while $\lambda_{I_i}$ also receives the
error carried by the voltage state.

At the arrival of spike $(j,k)$, the additive jump in $Q_i$ has an identity
state Jacobian. Hence all three adjoints are continuous across the arrival:
\begin{equation}
    \lambda_{V_i}^-=\lambda_{V_i}^+,
    \qquad
    \lambda_{I_i}^-=\lambda_{I_i}^+,
    \qquad
    \lambda_{Q_i}^-=\lambda_{Q_i}^+.
    \label{eq:tes-arrival-adjoint-continuity}
\end{equation}
The arrival generates the scalar timing signal
\begin{equation}
    \mu_{jk}
      =\frac{\partial l_{\mathrm p}}{\partial t^s_{jk}}
       +\sum_i w_{ij}
        \left(\omega_i\lambda_{I_i}^+
              +\alpha\lambda_{Q_i}^+\right)_{t^s_{jk}},
    \label{eq:tes-arrival-timing-signal}
\end{equation}
where the sum contains every postsynaptic target of neuron $j$. This scalar is
carried backward to the threshold event that scheduled $t^s_{jk}$.

For a threshold event $(i,k)$, define
\begin{align}
    v_{ik}&=\dot V^\theta_{ik},
    &a_{ik}&=\ddot V_i\!\left((t^\theta_{ik})^-\right),\\
    r_{ik}&=\sigma(c_1-c_0v_{ik}),
    &\mathfrak d'_{ik}&=-d_+c_0r_{ik}(1-r_{ik}).
    \label{eq:tes-threshold-auxiliary}
\end{align}
For the subtractive reset,
$\dot V_{ik}^{\theta,+}=v_{ik}+R_i/\tau_V$. The reverse threshold
updates are
\begin{align}
    \lambda_{V_i}^-
      &=\underbrace{
        \frac{\dot V_{ik}^{\theta,+}}{v_{ik}}\lambda_{V_i}^+
        +\frac{\mu_{ik}+l_V^- -l_V^+}{\tau_Vv_{ik}}
        }_{\text{EventProp-form term}}
        +\underbrace{
        \frac{\mu_{ik}\mathfrak d'_{ik}}{\tau_V}
        \left(\frac{1}{\tau_V}+\frac{a_{ik}}{v_{ik}}\right)
        }_{\text{state-dependent-delay correction}},
        \label{eq:tes-threshold-adjoint-v}\\
    \lambda_{I_i}^-
      &=\lambda_{I_i}^+
        -\frac{\mu_{ik}\mathfrak d'_{ik}}{\tau_V},
        \label{eq:tes-threshold-adjoint-i}\\
    \lambda_{Q_i}^-&=\lambda_{Q_i}^+.
        \label{eq:tes-threshold-adjoint-q}
\end{align}
The first two terms of Equation~\eqref{eq:tes-threshold-adjoint-v} have the
original EventProp structure. The final term and the jump in $\lambda_{I_i}$
are introduced by the dependence of the delay on the crossing slope.

Because a weight produces a negative jump of magnitude $w_{ij}$ in the
postsynaptic $Q$ state, its
gradient samples the corresponding adjoint at every transmitted spike:
\begin{equation}
    \frac{\mathrm d\mathcal L}{\mathrm d w_{ij}}
      =+\sum_{k:\,j\ \mathrm{spikes}}
        \lambda_{Q_i}\!\left(t^s_{jk}\right).
    \label{eq:tes-weight-gradient}
\end{equation}
When the oscillator frequencies are learned, introduce two secondary adjoint
helpers for every neuron, $\lambda_{I_i}^{\omega}$ and
$\lambda_{Q_i}^{\omega}$. They are propagated backward with terminal
conditions
\begin{equation}
    \lambda_{I_i}^{\omega}(T)
      =\lambda_{Q_i}^{\omega}(T)=0
\end{equation}
and continuous dynamics
\begin{align}
    (\lambda_{I_i}^{\omega})'
      &=-\alpha\lambda_{I_i}^{\omega}
        +\omega_i\lambda_{Q_i}^{\omega}
        -\lambda_{Q_i},
        \label{eq:tes-omega-helper-i}\\
    (\lambda_{Q_i}^{\omega})'
      &=-\omega_i\lambda_{I_i}^{\omega}
        -\alpha\lambda_{Q_i}^{\omega}
        +\lambda_{I_i}.
        \label{eq:tes-omega-helper-q}
\end{align}
Both helpers are continuous at thresholds and arrivals:
\begin{equation}
    (\lambda_{I_i}^{\omega})^-
      =(\lambda_{I_i}^{\omega})^+,
    \qquad
    (\lambda_{Q_i}^{\omega})^-
      =(\lambda_{Q_i}^{\omega})^+.
    \label{eq:tes-omega-helper-continuity}
\end{equation}
For zero initial synaptic state, the frequency gradient is then accumulated
only at arrivals into neuron $i$:
\begin{equation}
    \frac{\mathrm d\mathcal L}{\mathrm d\omega_i}
      =-\sum_{(j,k)\to i}
        w_{ij}\lambda_{Q_i}^{\omega}(t^s_{jk}).
    \label{eq:tes-omega-gradient}
\end{equation}
Thus frequency learning retains an event-sparse reverse pass without storing
the forward $I_i,Q_i$ trajectories.

\paragraph{Feed-forward continuous specialization used in this checkpoint.}
The implementation separates readout-timing objectives from neuron dynamics.
With target readout $y$, physical readout set $P$, and an objective-specific
margin $m$, define
\begin{equation}
    \tau_y=
    \begin{cases}
        t_y,&y\in P,\\
        t_a+2\,\mathrm{ms},&y\notin P,\quad
        a=\arg\max_{r\in P}t_r,
    \end{cases}
\end{equation}

The original DelGrad equality objective, used with
$m=1\,\mathrm{ms}$ in the 40-epoch comparison run, is
\begin{equation}
    \mathcal L_{\mathrm{DelGrad}}
      =\frac12\sum_{\substack{r\in P\\r\ne y}}
        \left((t_r-\tau_y)-m\right)^2.
    \label{eq:tes-masked-delgrad}
\end{equation}
The separately selectable squared-hinge objective uses
$m=2\,\mathrm{ms}$ by default:
\begin{equation}
    \mathcal L_{\mathrm{diag}}
      =\frac12\sum_{\substack{r\in P\\r\ne y}}
        \min\!\left\{(t_r-\tau_y)-m,0\right\}^2.
    \label{eq:tes-masked-squared-hinge}
\end{equation}
Under the hinge, a physical wrong readout exactly at or beyond the margin
contributes neither loss nor derivative. Under original DelGrad, every
physical wrong readout contributes. A fake wrong readout is excluded from
either objective. A fake
target is compared with every physical wrong readout, and the derivative of
its clipped time is routed through the unique latest physical anchor $a$.
The fake target itself still has zero timing source and no threshold event.
If the latest physical time is tied while the target is fake, the scalar loss
is reportable but its derivative is declared unavailable; no arbitrary
subgradient is used.  When all readouts are silent, the loss and
all timing sources are zero.

Because all first-spike events are already known in the feed-forward solver,
the reverse pass is performed layer by layer.  If
$\Phi_i(\Delta)=\exp(A_i\Delta)$ and
$M=\operatorname{diag}(\tau_V,1,1)$, exact reverse flow over an event-free
interval is
\begin{equation}
    \boldsymbol\lambda_i(t)
      =M^{-1}\Phi_i(\Delta)^\top M
       \boldsymbol\lambda_i(t+\Delta).
    \label{eq:tes-exact-layerwise-adjoint-flow}
\end{equation}
For a physical output $r$ with loss timing source
$g_r=\partial\mathcal L_{\mathrm{diag}}/\partial t_r$ and positive crossing
slope $v_r$, the pre-threshold seed is
\begin{equation}
    \boldsymbol\lambda_r^-
      =\begin{bmatrix}g_r/(\tau_Vv_r)&0&0\end{bmatrix}^{\!\top}.
\end{equation}
Reversing each hidden arrival $h\to r$ strictly before that threshold gives
\begin{align}
    \frac{\partial\mathcal L_{\mathrm{diag}}}
         {\partial W^{\mathrm{out}}_{rh}}
      &\mathrel{+}=\lambda_{Q_r},\\
    \mu_h
      &\mathrel{+}=W^{\mathrm{out}}_{rh}
        \left(\omega_r\lambda_{I_r}+\alpha\lambda_{Q_r}\right).
    \label{eq:tes-feedforward-readout-reverse}
\end{align}
All quantities on the right are sampled at the arrival time.  For each
physical hidden spike $h$, its accumulated $\mu_h$, slope $v_h$,
pre-threshold acceleration $a_h$, and delay derivative $\mathfrak d'_h$ enter
the full reset update
\begin{equation}
    \boldsymbol\lambda_h^-=
    \begin{bmatrix}
      \dfrac{v_h+R_h/\tau_V}{v_h}\lambda_{V_h}^+
      +\dfrac{\mu_h}{\tau_Vv_h}
      +\dfrac{\mu_h\mathfrak d'_h}{\tau_V}
       \left(\dfrac1{\tau_V}+\dfrac{a_h}{v_h}\right)\\[5pt]
      \lambda_{I_h}^+-\dfrac{\mu_h\mathfrak d'_h}{\tau_V}\\[5pt]
      \lambda_{Q_h}^+
    \end{bmatrix}.
    \label{eq:tes-feedforward-hidden-reset}
\end{equation}
Hidden thresholds, their emitted arrivals, and external arrivals are traversed
in descending time, so the post-reset adjoint in
Equation~\eqref{eq:tes-feedforward-hidden-reset} is generally nonzero after an
earlier threshold.  Reversing every exact external arrival before the latest
physical first-readout terminal then accumulates
\begin{equation}
    \frac{\partial\mathcal L_{\mathrm{diag}}}
         {\partial W^{\mathrm{in}}_{hc}}
      \mathrel{+}=\lambda_{Q_h}
\end{equation}
for every arrival on channel $c$.  Silent hidden rows and fake output rows
therefore remain exactly zero.  Oscillator frequencies are frozen in this
checkpoint.

For the proof-of-learning implementation, the exact equations above are retained
as the mathematical definition, while every occurrence of
$1/(\tau_V\dot V)$ in both continuous and grid-resolved reverse passes is
evaluated as $\min\{1/(\tau_V\dot V),10\}$. This numerical safeguard is a
cap-active approximation: gradients coincide with the exact derivation only
when the cap is inactive. Cap activations and the minimum causal crossing slope
are recorded separately from the optimizer's global gradient-norm clipping.
All initialization, training, validation, and test simulations use
\verb|early_stop=False|. Physical first spikes alone define prediction; silent
incorrect classes do not defeat a physical correct class, whereas all-silent
trials and exact earliest-time ties are incorrect. A hidden input row is
stochastically restarted only after causal silence on more than 500 of 1,000
validation examples: independent seeded unit-normal noise plus an offset of
$0.2$ is added to the four signed weights. The Adam moments and step are
deliberately retained. Training writes a checkpoint and an inspectable
per-epoch plot after every epoch, and the held-out test split is evaluated once
after the final validation-driven intervention.

\subsection{Derivation of the TeS backward equations}
\label{sec:eventprop-derivation}

We follow the adjoint derivation of EventProp~\cite{wunderlich2021eventprop}.
Fix one weight $w\equiv w_{ab}$, connecting presynaptic neuron $b$ to
postsynaptic neuron $a$. For a state variable $X$, write
$\partial_wX=\partial X/\partial w$ for its partial derivative at fixed time,
while $\mathrm dt_e/\mathrm dw$ denotes the total derivative of an event time.
We differentiate the augmented loss, integrate by parts, and then choose the
adjoint dynamics and jumps so that every occurrence of $\partial_wV$,
$\partial_wI$, and $\partial_wQ$ cancels.

\paragraph{Step 1: split the augmented loss at every event.}
Let
\begin{equation}
    0=t_0<t_1<\cdots<t_M<t_{M+1}=T
\end{equation}
contain every threshold time and every final spike time $t^s$ (including each
synaptic arrival). Between these events the forward equations can be written
in implicit form as
\begin{align}
    f_{V_i}&=\tau_V\dot V_i+V_i-I_i=0,\\
    f_{I_i}&=\dot I_i+\alpha I_i+\omega_iQ_i=0,\\
    f_{Q_i}&=\dot Q_i-\omega_iI_i+\alpha Q_i=0.
    \label{eq:tes-implicit-flow}
\end{align}
Adding these constraints with Lagrange multipliers on every event-free interval
gives
\begin{equation}
    \widetilde{\mathcal L}
      =l_{\mathrm p}\!\left(\{t^s_{jk}\}\right)
       +\sum_{m=0}^{M}\int_{t_m}^{t_{m+1}}
        \left[
          l_V
          +\sum_i\left(
             \lambda_{V_i}f_{V_i}
            +\lambda_{I_i}f_{I_i}
            +\lambda_{Q_i}f_{Q_i}
          \right)
        \right]\mathrm dt.
    \label{eq:tes-augmented-loss}
\end{equation}
Because $f_{V_i}=f_{I_i}=f_{Q_i}=0$ along the forward trajectory,
$\widetilde{\mathcal L}=\mathcal L$ and their total derivatives with respect
to $w$ are equal for any choice of the multipliers.

\paragraph{Step 2: differentiate the constraints with respect to the weight.}
The weight is constant in time, so $\partial_w$ commutes with
$\mathrm d/\mathrm dt$. Since $\alpha$, $\omega_i$, and $\tau_V$ are fixed,
\begin{align}
    \partial_w f_{V_i}
      &=\tau_V\frac{\mathrm d}{\mathrm dt}\partial_w V_i
        +\partial_w V_i-\partial_w I_i,
        \label{eq:tes-differentiate-fv}\\
    \partial_w f_{I_i}
      &=\frac{\mathrm d}{\mathrm dt}\partial_w I_i
        +\alpha\partial_w I_i+\omega_i\partial_w Q_i,
        \label{eq:tes-differentiate-fi}\\
    \partial_w f_{Q_i}
      &=\frac{\mathrm d}{\mathrm dt}\partial_w Q_i
        -\omega_i\partial_w I_i+\alpha\partial_w Q_i.
        \label{eq:tes-differentiate-fq}
\end{align}

\paragraph{Step 3: apply the Leibniz rule.}
Differentiating Equation~\eqref{eq:tes-augmented-loss}, substituting
Equations~\eqref{eq:tes-differentiate-fv}--
\eqref{eq:tes-differentiate-fq}, and using $f_{V_i}=f_{I_i}=f_{Q_i}=0$
gives the expression below. In particular,
$\partial_w(\lambda_X f_X)
=(\partial_w\lambda_X)f_X+\lambda_X\partial_wf_X$ reduces to
$\lambda_X\partial_wf_X$ because $f_X=0$. At the moving interval boundaries
the same constraint terms vanish, leaving only the $l_V$ boundary values:
\begin{align}
    \frac{\mathrm d\mathcal L}{\mathrm dw}
    ={}&\sum_{j,k}
       \frac{\partial l_{\mathrm p}}{\partial t^s_{jk}}
       \frac{\mathrm dt^s_{jk}}{\mathrm dw}
       \nonumber\\
    &+\sum_{m=0}^{M}\int_{t_m}^{t_{m+1}}
      \sum_i\Bigg[
       \frac{\partial l_V}{\partial V_i}\partial_w V_i
       +\lambda_{V_i}\left(
          \tau_V\frac{\mathrm d}{\mathrm dt}\partial_w V_i
          +\partial_w V_i-\partial_w I_i\right)
       \nonumber\\
    &\hspace{31mm}
       +\lambda_{I_i}\left(
          \frac{\mathrm d}{\mathrm dt}\partial_w I_i
          +\alpha\partial_w I_i+\omega_i\partial_w Q_i\right)
       \nonumber\\
    &\hspace{31mm}
       +\lambda_{Q_i}\left(
          \frac{\mathrm d}{\mathrm dt}\partial_w Q_i
          -\omega_i\partial_w I_i+\alpha\partial_w Q_i\right)
      \Bigg]\mathrm dt
      \nonumber\\
    &+\sum_{m=0}^{M}\left[
       l_{V,m+1}^-\frac{\mathrm dt_{m+1}}{\mathrm dw}
       -l_{V,m}^+\frac{\mathrm dt_m}{\mathrm dw}
      \right].
    \label{eq:tes-gradient-before-parts}
\end{align}
Here $l_{V,m}^{-}$ and $l_{V,m}^{+}$ are the running loss immediately before
and after event $t_m$. The constraint terms do not contribute at the moving
integration limits because every $f$ is zero.

\paragraph{Step 4: integrate all time derivatives by parts.}
On each event-free interval,
\begin{align}
    \int_{t_m}^{t_{m+1}}
      \tau_V\lambda_{V_i}
      \frac{\mathrm d}{\mathrm dt}\partial_w V_i\,\mathrm dt
      &=
      \tau_V\left[
        \lambda_{V_i}\partial_w V_i
      \right]_{t_m^+}^{t_{m+1}^-}
      -\int_{t_m}^{t_{m+1}}
        \tau_V\dot\lambda_{V_i}\partial_w V_i\,\mathrm dt,
        \label{eq:tes-parts-v}\\
    \int_{t_m}^{t_{m+1}}
      \lambda_{I_i}
      \frac{\mathrm d}{\mathrm dt}\partial_w I_i\,\mathrm dt
      &=
      \left[
        \lambda_{I_i}\partial_w I_i
      \right]_{t_m^+}^{t_{m+1}^-}
      -\int_{t_m}^{t_{m+1}}
        \dot\lambda_{I_i}\partial_w I_i\,\mathrm dt,
        \label{eq:tes-parts-i}\\
    \int_{t_m}^{t_{m+1}}
      \lambda_{Q_i}
      \frac{\mathrm d}{\mathrm dt}\partial_w Q_i\,\mathrm dt
      &=
      \left[
        \lambda_{Q_i}\partial_w Q_i
      \right]_{t_m^+}^{t_{m+1}^-}
      -\int_{t_m}^{t_{m+1}}
        \dot\lambda_{Q_i}\partial_w Q_i\,\mathrm dt.
        \label{eq:tes-parts-q}
\end{align}
Substitution into Equation~\eqref{eq:tes-gradient-before-parts} and collection
of equal state derivatives gives
\begin{align}
    \frac{\mathrm d\mathcal L}{\mathrm dw}
    ={}&\sum_{m=0}^{M}\int_{t_m}^{t_{m+1}}\sum_i
      \Bigg[
      \underbrace{\left(
        \frac{\partial l_V}{\partial V_i}
        +\lambda_{V_i}-\tau_V\dot\lambda_{V_i}
      \right)}_{C_{V_i}}\partial_w V_i
      \nonumber\\
    &\hspace{21mm}
      +\underbrace{\left(
        -\lambda_{V_i}+\alpha\lambda_{I_i}
        -\omega_i\lambda_{Q_i}-\dot\lambda_{I_i}
      \right)}_{C_{I_i}}\partial_w I_i
      \nonumber\\
    &\hspace{21mm}
      +\underbrace{\left(
        \omega_i\lambda_{I_i}
        +\alpha\lambda_{Q_i}-\dot\lambda_{Q_i}
      \right)}_{C_{Q_i}}\partial_w Q_i
      \Bigg]\mathrm dt
      +B,
    \label{eq:tes-gradient-collected}
\end{align}
where no term has been discarded: the complete remainder is
\begin{align}
    B={}&
      \sum_{j,k}
       \frac{\partial l_{\mathrm p}}{\partial t^s_{jk}}
       \frac{\mathrm dt^s_{jk}}{\mathrm dw}
      +\sum_{m=0}^{M}\left[
       l_{V,m+1}^-\frac{\mathrm dt_{m+1}}{\mathrm dw}
       -l_{V,m}^+\frac{\mathrm dt_m}{\mathrm dw}
      \right]
      \nonumber\\
    &+\sum_{m=0}^{M}\sum_i\Bigg\{
      \tau_V\left[
        \lambda_{V_i}\partial_wV_i
      \right]_{t_m^+}^{t_{m+1}^-}
      +\left[
        \lambda_{I_i}\partial_wI_i
      \right]_{t_m^+}^{t_{m+1}^-}
      \nonumber\\
    &\hspace{34mm}
      +\left[
        \lambda_{Q_i}\partial_wQ_i
      \right]_{t_m^+}^{t_{m+1}^-}
      \Bigg\}.
    \label{eq:tes-complete-boundary-remainder}
\end{align}

\paragraph{Step 5: choose the continuous adjoint equations.}
The multipliers are free because the constraints vanish. We choose them so
that the three coefficients above vanish separately:
\begin{align}
    C_{V_i}=0
      &\quad\Longrightarrow\quad
      \tau_V\dot\lambda_{V_i}
        =\lambda_{V_i}+\frac{\partial l_V}{\partial V_i},
        \label{eq:tes-forward-time-adjoint-v}\\
    C_{I_i}=0
      &\quad\Longrightarrow\quad
      \dot\lambda_{I_i}
        =-\lambda_{V_i}+\alpha\lambda_{I_i}
         -\omega_i\lambda_{Q_i},
        \label{eq:tes-forward-time-adjoint-i}\\
    C_{Q_i}=0
      &\quad\Longrightarrow\quad
      \dot\lambda_{Q_i}
        =\omega_i\lambda_{I_i}+\alpha\lambda_{Q_i}.
        \label{eq:tes-forward-time-adjoint-q}
\end{align}
This makes every interval integral in
Equation~\eqref{eq:tes-gradient-collected} exactly zero. Replacing
$\mathrm d/\mathrm dt$ by $-\mathrm d/\mathrm dt$ for reverse-time
integration gives Equations~\eqref{eq:tes-adjoint-v}--
\eqref{eq:tes-adjoint-q}.

\paragraph{Step 6: rearrange the interval endpoints into event terms.}
For example, the voltage endpoints satisfy
\begin{align}
    &\sum_{m=0}^{M}
      \tau_V\left[
        \lambda_{V_i}\partial_wV_i
      \right]_{t_m^+}^{t_{m+1}^-}
      \nonumber\\
    &\quad=
      \tau_V\lambda_{V_i}(T)\partial_wV_i(T)
      -\tau_V\lambda_{V_i}(0)\partial_wV_i(0)
      \nonumber\\
    &\qquad+
      \sum_{m=1}^{M}\tau_V\left(
        \lambda_{V_i}^-\partial_wV_i^-
        -\lambda_{V_i}^+\partial_wV_i^+
      \right)_{t_m}.
    \label{eq:tes-telescope-v}
\end{align}
The $I_i$ and $Q_i$ endpoint sums rearrange identically.
Parameter-independent initial states give
$\partial_wV_i(0)=\partial_wI_i(0)=\partial_wQ_i(0)=0$, and we impose
$\lambda_{V_i}(T)=\lambda_{I_i}(T)=\lambda_{Q_i}(T)=0$. Hence the initial
and terminal terms vanish. The moving-limit terms telescope as
\(
\mathrm dt_0/\mathrm dw
=\mathrm dt_{M+1}/\mathrm dw=0
\), so
\begin{equation}
    \sum_{m=0}^{M}\left[
       l_{V,m+1}^-\frac{\mathrm dt_{m+1}}{\mathrm dw}
       -l_{V,m}^+\frac{\mathrm dt_m}{\mathrm dw}
    \right]
    =
    \sum_{m=1}^{M}
      (l_{V,m}^- -l_{V,m}^+)
      \frac{\mathrm dt_m}{\mathrm dw}.
    \label{eq:tes-telescope-loss}
\end{equation}
After these rearrangements,
\begin{equation}
    \frac{\mathrm d\mathcal L}{\mathrm dw}
      =\sum_{m=1}^{M}\xi_m,
    \label{eq:tes-gradient-event-sum}
\end{equation}
with
\begin{align}
    \xi_m={}&
      \frac{\partial l_{\mathrm p}}{\partial t^s_{m}}
      \frac{\mathrm dt^s_{m}}{\mathrm dw}
      +(l_{V,m}^- -l_{V,m}^+)\frac{\mathrm dt_m}{\mathrm dw}
      \nonumber\\
    &+\sum_i\Bigg[
      \tau_V\left(
        \lambda_{V_i}^-\partial_wV_i^-
        -\lambda_{V_i}^+\partial_wV_i^+
      \right)
      \nonumber\\
    &\hspace{17mm}
      +\lambda_{I_i}^-\partial_wI_i^-
       -\lambda_{I_i}^+\partial_wI_i^+
      +\lambda_{Q_i}^-\partial_wQ_i^-
       -\lambda_{Q_i}^+\partial_wQ_i^+
      \Bigg]_{t_m}.
    \label{eq:tes-event-contribution}
\end{align}
This is the EventProp boundary sum specialized to the three TeS states.
We now differentiate each transition equality directly and substitute the
result into Equation~\eqref{eq:tes-event-contribution}.

\paragraph{Step 7: partial derivatives across a post-synaptic arrival.}
Consider the arrival of spike $(j,k)$ at a postsynaptic neuron $i$ and
differentiate with respect to $w=w_{ab}$. The forward transition is
\begin{equation}
    V_i^+=V_i^-,
    \qquad
    I_i^+=I_i^-,
    \qquad
    Q_i^+=Q_i^- -w_{ij}.
    \label{eq:tes-arrival-transition-proof}
\end{equation}
Substitution of these jumps into the forward differential equations gives
\begin{align}
    \dot V_i^+-\dot V_i^-
      &=0,
      \nonumber\\
    \dot I_i^+-\dot I_i^-
      &
        =\omega_iw_{ij},
      \nonumber\\
    \dot Q_i^+-\dot Q_i^-
      &
        =\alpha w_{ij}.
    \label{eq:tes-arrival-flow-change}
\end{align}

Equation~\eqref{eq:tes-arrival-transition-proof} holds at the weight-dependent
time $t^s_{jk}$. Differentiating each equality with respect to $w$ therefore
gives
\begin{align}
    \partial_wV_i^+
      +\dot V_i^+\frac{\mathrm dt^s_{jk}}{\mathrm dw}
      &=
      \partial_wV_i^-
      +\dot V_i^-\frac{\mathrm dt^s_{jk}}{\mathrm dw},
      \label{eq:tes-arrival-differentiate-v}\\
    \partial_wI_i^+
      +\dot I_i^+\frac{\mathrm dt^s_{jk}}{\mathrm dw}
      &=
      \partial_wI_i^-
      +\dot I_i^-\frac{\mathrm dt^s_{jk}}{\mathrm dw},
      \label{eq:tes-arrival-differentiate-i}\\
    \partial_wQ_i^+
      +\dot Q_i^+\frac{\mathrm dt^s_{jk}}{\mathrm dw}
      &=
      \partial_wQ_i^-
      +\dot Q_i^-\frac{\mathrm dt^s_{jk}}{\mathrm dw}
      -\delta_{ia}\delta_{jb}.
      \label{eq:tes-arrival-differentiate-q}
\end{align}
Rearranging these three equalities and using
Equation~\eqref{eq:tes-arrival-flow-change} yields
\begin{align}
    \partial_wV_i^+
      &=\partial_wV_i^-,
      \nonumber\\
    \partial_wI_i^+
      &=\partial_wI_i^-
        +(\dot I_i^- -\dot I_i^+)
          \frac{\mathrm dt^s_{jk}}{\mathrm dw}
       =\partial_wI_i^-
        -\omega_iw_{ij}\frac{\mathrm dt^s_{jk}}{\mathrm dw},
      \nonumber\\
    \partial_wQ_i^+
      &=\partial_wQ_i^-
        +(\dot Q_i^- -\dot Q_i^+)
          \frac{\mathrm dt^s_{jk}}{\mathrm dw}
        -\delta_{ia}\delta_{jb}
      \nonumber\\
      &=\partial_wQ_i^-
        -\alpha w_{ij}\frac{\mathrm dt^s_{jk}}{\mathrm dw}
        -\delta_{ia}\delta_{jb}.
    \label{eq:tes-arrival-partial-derivatives}
\end{align}
The Kronecker term is present only when the arrival crosses the differentiated
synapse. Because $V_i$ is continuous, $l_V^-=l_V^+$ at an arrival.
Substituting Equation~\eqref{eq:tes-arrival-partial-derivatives} into
Equation~\eqref{eq:tes-event-contribution}, and summing over all postsynaptic
targets of spike $(j,k)$, gives
\begin{align}
    \xi^s_{jk}
    ={}&\sum_i\Big[
       \tau_V(\lambda_{V_i}^- -\lambda_{V_i}^+)\partial_wV_i^-
       +(\lambda_{I_i}^- -\lambda_{I_i}^+)\partial_wI_i^-
       \nonumber\\
    &\hspace{18mm}
       +(\lambda_{Q_i}^- -\lambda_{Q_i}^+)\partial_wQ_i^-
       +\lambda_{Q_i}^+\delta_{ia}\delta_{jb}
       \Big]
       \nonumber\\
    &+\left[
       \frac{\partial l_{\mathrm p}}{\partial t^s_{jk}}
       +\sum_iw_{ij}
        (\omega_i\lambda_{I_i}^+
         +\alpha\lambda_{Q_i}^+)
       \right]\frac{\mathrm dt^s_{jk}}{\mathrm dw}.
    \label{eq:tes-arrival-boundary-expansion}
\end{align}
We do not want to calculate the three state partial derivatives. Their
coefficients vanish when the arrival adjoints are chosen to be continuous:
\begin{equation}
    \lambda_{V_i}^-=\lambda_{V_i}^+,
    \qquad
    \lambda_{I_i}^-=\lambda_{I_i}^+,
    \qquad
    \lambda_{Q_i}^-=\lambda_{Q_i}^+.
    \label{eq:tes-arrival-adjoint-continuity-derived}
\end{equation}
This is Equation~\eqref{eq:tes-arrival-adjoint-continuity}. Define the
remaining coefficient of $\mathrm dt^s_{jk}/\mathrm dw$ as
\begin{equation}
    \mu_{jk}
      =\frac{\partial l_{\mathrm p}}{\partial t^s_{jk}}
       +\sum_i w_{ij}
        \left(\omega_i\lambda_{I_i}^+
              +\alpha\lambda_{Q_i}^+\right)_{t^s_{jk}},
    \label{eq:tes-arrival-timing-signal-derived}
\end{equation}
which is Equation~\eqref{eq:tes-arrival-timing-signal}. The arrival
contribution has now been reduced to
\begin{equation}
    \xi^s_{jk}
      =\mu_{jk}\frac{\mathrm dt^s_{jk}}{\mathrm dw}
       +\sum_i\lambda_{Q_i}^+
        \delta_{ia}\delta_{jb}.
    \label{eq:tes-arrival-after-cancellation}
\end{equation}
There is no direct $\lambda_V$ term in $\mu$, because $V_i$, $I_i$, and
$\dot V_i$ do not jump at a $Q$-injected arrival. Voltage error enters
indirectly through the $\lambda_V$ term in the continuous equation for
$\lambda_I$.

For external input spikes, $\frac{\mathrm dt^s_{\mathrm{in}}}{\mathrm dw} = 0$, thus only the Kronecker term remain nonzero.

\paragraph{Step 8: partial derivatives at the pre-synaptic threshold crossing.}
Now consider the threshold event that emitted the delayed spike from neuron $j$ and suppress the indices $(j,k)$. There is no any jumping response at post-synaptic neurons, thus their terms are all cancelled out, leaving only pre-synaptic terms. Write $v=\dot V((t^\theta)^-)>0$ and
$a=\ddot V((t^\theta)^-)$. The threshold equation is
\begin{equation}
    V^-(t^\theta)-\vartheta=0.
\end{equation}
Differentiating this equality with respect to $w$ gives
\begin{equation}
    \partial_wV^-+v\frac{\mathrm dt^\theta}{\mathrm dw}=0,
    \label{eq:tes-threshold-differentiated}
\end{equation}
and therefore
\begin{equation}
    \frac{\mathrm dt^\theta}{\mathrm dw}
      =-\frac{1}{v}\partial_wV^-.
    \label{eq:tes-threshold-time-variation}
\end{equation}

The crossing slope is $v=\dot V^-$. Its total derivative contains the
fixed-time partial derivative of $\dot V^-$ and the derivative caused by the change of threshold time:
\begin{equation}
    \frac{\mathrm dv}{\mathrm dw}
      =\partial_w\dot V^-
       +a\frac{\mathrm dt^\theta}{\mathrm dw}.
    \label{eq:tes-slope-chain-rule}
\end{equation}
From $\dot V=(-V+I)/\tau_V$,
\begin{equation}
    \partial_w\dot V^-
      =-\frac{1}{\tau_V}\partial_wV^-
       +\frac{1}{\tau_V}\partial_wI^-.
    \label{eq:tes-slope-fixed-time-partial}
\end{equation}
Substituting Equations~\eqref{eq:tes-threshold-time-variation} and
\eqref{eq:tes-slope-fixed-time-partial} into
Equation~\eqref{eq:tes-slope-chain-rule} gives
\begin{equation}
    \frac{\mathrm dv}{\mathrm dw}
      =-\left(\frac{1}{\tau_V}+\frac{a}{v}\right)\partial_wV^-
       +\frac{1}{\tau_V}\partial_wI^-.
    \label{eq:tes-slope-variation}
\end{equation}
The final spike time satisfies $t^s=t^\theta+\mathfrak d(v)$. Hence
\begin{align}
    \frac{\mathrm dt^s}{\mathrm dw}
      &=\frac{\mathrm dt^\theta}{\mathrm dw}
        +\mathfrak d'\frac{\mathrm dv}{\mathrm dw}
        \nonumber\\
      &=-\left[
          \frac{1}{v}
          +\mathfrak d'
           \left(\frac{1}{\tau_V}+\frac{a}{v}\right)
        \right]\partial_wV^-
       +\frac{\mathfrak d'}{\tau_V}\partial_wI^-.
    \label{eq:tes-final-time-variation}
\end{align}

The subtractive reset is defined by
\begin{equation}
    V^+=V^- -R,
    \qquad I^+=I^-,
    \qquad Q^+=Q^-.
    \label{eq:tes-threshold-transition-proof}
\end{equation}
Differentiating these three equalities at $t^\theta(w)$ gives
\begin{align}
    \partial_wV^+
      +\dot V^{\theta,+}\frac{\mathrm dt^\theta}{\mathrm dw}
      &=
      \partial_wV^-
      +v\frac{\mathrm dt^\theta}{\mathrm dw},
      \label{eq:tes-reset-differentiate-v}\\
    \partial_wI^+
      +\dot I^+\frac{\mathrm dt^\theta}{\mathrm dw}
      &=
      \partial_wI^-
      +\dot I^-\frac{\mathrm dt^\theta}{\mathrm dw},
      \label{eq:tes-reset-differentiate-i}\\
    \partial_wQ^+
      +\dot Q^+\frac{\mathrm dt^\theta}{\mathrm dw}
      &=
      \partial_wQ^-
      +\dot Q^-\frac{\mathrm dt^\theta}{\mathrm dw}.
      \label{eq:tes-reset-differentiate-q}
\end{align}
At a threshold reset, $I$ and $Q$ do not jump, so their vector fields also do
not jump. Using Equation~\eqref{eq:tes-threshold-time-variation} in the first
line and $\dot I^+=\dot I^-$, $\dot Q^+=\dot Q^-$ in the other two lines gives
\begin{equation}
    \partial_wV^+
      =\frac{\dot V^{\theta,+}}{v}\partial_wV^-,
    \qquad
    \partial_wI^+=\partial_wI^-,
    \qquad
    \partial_wQ^+=\partial_wQ^-.
    \label{eq:tes-reset-partial-derivatives}
\end{equation}
For the subtractive reset,
$\dot V^{\theta,+}=v+R/\tau_V$. Let $\lambda_Q^-=\lambda_Q^+$, we get:
\begin{align}
    \xi^\theta={}&
     \left(l_V^- -l_V^+\right)\frac{\mathrm dt^\theta}{\mathrm dw}
     +\tau_V\left(
        \lambda_V^-\partial_w V_j^-
        -\lambda_V^+\partial_w V_j^+\right) +(\lambda_{I_j}^-
       -\lambda_{I_j}^+)\partial_wI_j^-
    \label{eq:tes-threshold_pre_threshold}
\end{align}

\paragraph{Step 9: cancel the threshold partial derivatives.}
Each pre-synaptic crossing is corresponding to a post-synaptic arrival, whose equation~\eqref{eq:tes-arrival-after-cancellation} leaves $\mu\,\mathrm dt^s/\mathrm dw$ for this emitted spike. We combine pre- and post-synaptic (shift $\mathrm dt^s$ to $\mathrm dt^\theta$) terms from
Equation~\eqref{eq:tes-event-contribution}:
\begin{align}
    \xi^\theta_{jk} + \xi^s_{jk}={}&
     \left(l_V^- -l_V^+\right)\frac{\mathrm dt^\theta}{\mathrm dw}
     +\tau_V\left(
        \lambda_V^-\partial_wV^-
        -\lambda_V^+\partial_wV^+\right)
     +(\lambda_{I_j}^-
       -\lambda_{I_j}^+)\partial_wI_j^-
     \nonumber\\
    &
     +\mu_{jk}\left(\left[
          \frac{1}{v}
          +\mathfrak d'
           \left(\frac{1}{\tau_V}+\frac{a}{v}\right)
        \right]\partial_wV^-
       +\frac{\mathfrak d'}{\tau_V}\partial_wI^-\right).
    \label{eq:tes-threshold-boundary-before-substitution}
\end{align}
Substitution of Equations~\eqref{eq:tes-threshold-time-variation}--
\eqref{eq:tes-reset-partial-derivatives} yields
\begin{align}
    \xi^\theta={}&
      \Bigg\{
        \tau_V\left(
        \lambda_V^-
        -\frac{\dot V^{\theta,+}}{v}\lambda_V^+
        \right)
        -\frac{l_V^- -l_V^+}{v}
        -\mu\left[
          \frac{1}{v}
          +\mathfrak d'\left(\frac{1}{\tau_V}+\frac{a}{v}\right)
        \right]
      \Bigg\}\partial_wV^-
      \nonumber\\
    &+\left(
        \lambda_I^- -\lambda_I^+
        +\frac{\mu\mathfrak d'}{\tau_V}
      \right)\partial_wI^-
    \label{eq:tes-threshold-collected}
\end{align}
We choose the threshold adjoint jumps so that each of the three coefficients
vanishes. 
The $I$ coefficient gives
\begin{equation}
    \lambda_I^- -\lambda_I^+
      +\frac{\mu\mathfrak d'}{\tau_V}=0
    \quad\Longrightarrow\quad
    \lambda_I^-=\lambda_I^+
      -\frac{\mu\mathfrak d'}{\tau_V}.
    \label{eq:tes-boundary-coefficient-i}
\end{equation}
Finally, setting the coefficient of $\partial_wV^-$ to zero and isolating
$\lambda_V^-$ gives
\begin{align}
    \tau_V\lambda_V^-
      ={}&\tau_V\frac{\dot V^{\theta,+}}{v}\lambda_V^+
       +\frac{l_V^- -l_V^+}{v}
       \nonumber\\
      &+\mu\left[
          \frac{1}{v}
          +\mathfrak d'\left(
             \frac{1}{\tau_V}+\frac{a}{v}
          \right)
       \right],
    \label{eq:tes-boundary-coefficient-v-before-division}\\
    \lambda_V^-
      ={}&\frac{\dot V^{\theta,+}}{v}\lambda_V^+
       +\frac{\mu+l_V^- -l_V^+}{\tau_Vv}
       +\frac{\mu\mathfrak d'}{\tau_V}
        \left(\frac{1}{\tau_V}+\frac{a}{v}\right).
    \label{eq:tes-boundary-coefficient-v}
\end{align}
These are precisely Equations~\eqref{eq:tes-threshold-adjoint-v}--
\eqref{eq:tes-threshold-adjoint-q}. With these choices,
$\xi^\theta + \xi^s=0$: both the threshold boundary terms and the carried $\mu\,\mathrm dt^s/\mathrm dw$ term cancel without computing any state partial derivative. For readout spikes without post-synaptic target, only \eqref{eq:tes-threshold_pre_threshold} need to be taken care.

\paragraph{Step 10: retain only the direct weight term.}
The interval terms vanished in Step~5, the arrival state derivatives vanished
in Step~7, and the threshold terms vanished in Step~9. Therefore only the
Kronecker term in Equation~\eqref{eq:tes-arrival-after-cancellation} remains:
\begin{align}
    \frac{\mathrm d\mathcal L}{\mathrm dw_{ab}}
      &=\sum_{j,k}\sum_i
        \lambda_{Q_i}(t^s_{jk})
        \delta_{ia}\delta_{jb}
        \nonumber\\
      &=\sum_{k:\,b\ \mathrm{spikes}}
        \lambda_{Q_a}(t^s_{bk}),
    \label{eq:tes-weight-gradient-derived}
\end{align}
which is Equation~\eqref{eq:tes-weight-gradient}.

\subsubsection{Oscillator-frequency gradient}
\label{sec:tes-omega-gradient-derivation}

We next differentiate the same augmented loss,
Equation~\eqref{eq:tes-augmented-loss}, with respect to one oscillator
frequency $\omega_a$. This first adjoint stage follows the weight derivation
above. Its event-boundary terms have the same form as those for a weight; the
only new interval terms arise because $\omega_a$ occurs explicitly in the two
oscillator constraints.

\paragraph{Step 11: differentiate the same augmented loss with respect to
$\omega_a$.}
At fixed time, the differentiated constraints are
\begin{align}
    \partial_{\omega_a}f_{V_i}
      &=\tau_V\frac{\mathrm d}{\mathrm dt}
          \partial_{\omega_a}V_i
        +\partial_{\omega_a}V_i
        -\partial_{\omega_a}I_i,
        \label{eq:tes-omega-differentiate-fv}\\
    \partial_{\omega_a}f_{I_i}
      &=\frac{\mathrm d}{\mathrm dt}\partial_{\omega_a}I_i
        +\alpha\partial_{\omega_a}I_i
        +\omega_i\partial_{\omega_a}Q_i
        +\delta_{ia}Q_i,
        \label{eq:tes-omega-differentiate-fi}\\
    \partial_{\omega_a}f_{Q_i}
      &=\frac{\mathrm d}{\mathrm dt}\partial_{\omega_a}Q_i
        -\omega_i\partial_{\omega_a}I_i
        +\alpha\partial_{\omega_a}Q_i
        -\delta_{ia}I_i.
        \label{eq:tes-omega-differentiate-fq}
\end{align}
The final terms in the last two equations are the explicit derivatives of
$\omega_iQ_i$ and $-\omega_iI_i$, respectively.

Applying the same Leibniz rule and integration by parts as in
Steps~3--4 gives the same coefficients $C_{V_i}$, $C_{I_i}$, and $C_{Q_i}$
as Equation~\eqref{eq:tes-gradient-collected}, together with the two explicit
frequency terms:
\begin{align}
    \frac{\mathrm d\mathcal L}{\mathrm d\omega_a}
    ={}&\sum_{m=0}^{M}\int_{t_m}^{t_{m+1}}\sum_i
      \Big[
        C_{V_i}\partial_{\omega_a}V_i
       +C_{I_i}\partial_{\omega_a}I_i
       +C_{Q_i}\partial_{\omega_a}Q_i
      \Big]\,\mathrm dt
      \nonumber\\
    &+\int_0^T
      \left(\lambda_{I_a}Q_a-\lambda_{Q_a}I_a\right)\,\mathrm dt
      +B_{\omega_a}.
    \label{eq:tes-omega-gradient-before-cancellation}
\end{align}
Here $B_{\omega_a}$ is exactly the complete boundary remainder in
Equation~\eqref{eq:tes-complete-boundary-remainder} after replacing every
$\partial_wX$ and $\mathrm dt_e/\mathrm dw$ by
$\partial_{\omega_a}X$ and $\mathrm dt_e/\mathrm d\omega_a$.
Choosing the ordinary adjoints according to
Equations~\eqref{eq:tes-forward-time-adjoint-v}--
\eqref{eq:tes-forward-time-adjoint-q} again sets all three $C$ coefficients
to zero. Therefore,
\begin{equation}
    \frac{\mathrm d\mathcal L}{\mathrm d\omega_a}
      =\int_0^T
        \left(\lambda_{I_a}Q_a-\lambda_{Q_a}I_a\right)\,\mathrm dt
       +\sum_{m=1}^{M}\xi_m^{\omega_a},
    \label{eq:tes-omega-gradient-event-remainder}
\end{equation}
where $\xi_m^{\omega_a}$ is the same EventProp boundary expression as
Equation~\eqref{eq:tes-event-contribution}, with the differentiated parameter
changed from $w$ to $\omega_a$.

\paragraph{Step 12: show that the frequency event remainder vanishes.}
At an arrival $t^s_{jk}$, the transition remains
$Q_i^+=Q_i^- -w_{ij}$. The weight $w_{ij}$ is independent of $\omega_a$.
Consequently, the frequency analogue of
Equation~\eqref{eq:tes-arrival-differentiate-q} is
\begin{equation}
    \partial_{\omega_a}Q_i^+
      +\dot Q_i^+\frac{\mathrm dt^s_{jk}}{\mathrm d\omega_a}
      =
      \partial_{\omega_a}Q_i^-
      +\dot Q_i^-\frac{\mathrm dt^s_{jk}}{\mathrm d\omega_a},
    \label{eq:tes-omega-arrival-differentiate-q}
\end{equation}
with no Kronecker term. The $V_i$ and $I_i$ equalities are obtained from
Equations~\eqref{eq:tes-arrival-differentiate-v} and
\eqref{eq:tes-arrival-differentiate-i} by the same replacement
$\partial_w\mapsto\partial_{\omega_a}$.
Using Equation~\eqref{eq:tes-arrival-flow-change}, the three fixed-time
partials satisfy
\begin{align}
    \partial_{\omega_a}V_i^+
      &=\partial_{\omega_a}V_i^-,
      \nonumber\\
    \partial_{\omega_a}I_i^+
      &=\partial_{\omega_a}I_i^-
        -\omega_iw_{ij}
          \frac{\mathrm dt^s_{jk}}{\mathrm d\omega_a},
      \nonumber\\
    \partial_{\omega_a}Q_i^+
      &=\partial_{\omega_a}Q_i^-
        -\alpha w_{ij}
          \frac{\mathrm dt^s_{jk}}{\mathrm d\omega_a}.
    \label{eq:tes-omega-arrival-partial-derivatives}
\end{align}
Substitution into the event boundary expression and use of the ordinary
arrival-adjoint continuity conditions gives
\begin{equation}
    \xi_{jk}^{s,\omega_a}
      =\mu_{jk}\frac{\mathrm dt^s_{jk}}{\mathrm d\omega_a}.
    \label{eq:tes-omega-arrival-after-cancellation}
\end{equation}
In contrast to Equation~\eqref{eq:tes-arrival-after-cancellation}, there is
no direct $\lambda_{Q_i}\delta_{ia}\delta_{jb}$ contribution. This is the
only difference between the two first-stage event calculations. The scalar
$\mu_{jk}$ is the boundary coefficient produced by the adjoint rearrangement
and is evaluated on the forward trajectory; it is not differentiated again
with respect to $\omega_a$.

At the presynaptic threshold, Equations~\eqref{eq:tes-threshold-differentiated}--
\eqref{eq:tes-reset-partial-derivatives} remain valid after replacing
$w$ by $\omega_a$: the soma equation, reset map, and delay map have no
explicit dependence on $\omega_a$ at fixed state. The same threshold adjoint
jumps, Equations~\eqref{eq:tes-threshold-adjoint-v}--
\eqref{eq:tes-threshold-adjoint-q}, therefore cancel
$\mu_{jk}\,\mathrm dt^s_{jk}/\mathrm d\omega_a$ exactly. Thresholds do not
leave a direct frequency term either. Hence
\begin{equation}
    \sum_{m=1}^{M}\xi_m^{\omega_a}=0,
\end{equation}
and the first adjoint stage gives
\begin{equation}
    \frac{\mathrm d\mathcal L}{\mathrm d\omega_i}
      =\int_0^T
        \left(\lambda_{I_i}Q_i-\lambda_{Q_i}I_i\right)\,\mathrm dt.
    \label{eq:tes-omega-gradient-integral}
\end{equation}

\paragraph{Step 13: introduce the second adjoint augmentation.}
Equation~\eqref{eq:tes-omega-gradient-integral} is already the complete
frequency derivative. We now rewrite that derivative; we do not differentiate
it again. Define
\begin{equation}
    \mathcal G_i
      =\int_0^T
        \left(\lambda_{I_i}Q_i-\lambda_{Q_i}I_i\right)\,\mathrm dt
      =\frac{\mathrm d\mathcal L}{\mathrm d\omega_i}.
\end{equation}
To remove the need to retain the forward $I_i,Q_i$ trajectory, augment
$\mathcal G_i$ on every event-free interval with the same oscillator
constraints, now using the secondary adjoints
$\lambda_{I_i}^{\omega}$ and $\lambda_{Q_i}^{\omega}$:
\begin{align}
    \widetilde{\mathcal G}_i
      =\sum_{m=0}^{M}\int_{t_m}^{t_{m+1}}
       \Big[
        &\lambda_{I_i}Q_i-\lambda_{Q_i}I_i
        -\lambda_{I_i}^{\omega}
          (\dot I_i+\alpha I_i+\omega_iQ_i)
        \nonumber\\
        &-\lambda_{Q_i}^{\omega}
          (\dot Q_i-\omega_iI_i+\alpha Q_i)
       \Big]\,\mathrm dt.
    \label{eq:tes-omega-second-augmented-functional}
\end{align}
The minus signs fix the orientation of the reverse helper equations. Because
the two forward constraints vanish,
$\widetilde{\mathcal G}_i=\mathcal G_i$ for arbitrary helpers.

Integrating the terms containing $\dot I_i$ and $\dot Q_i$ by parts gives
\begin{align}
    -\int\lambda_{I_i}^{\omega}\dot I_i\,\mathrm dt
      &=-[\lambda_{I_i}^{\omega}I_i]
        +\int\dot\lambda_{I_i}^{\omega}I_i\,\mathrm dt,\\
    -\int\lambda_{Q_i}^{\omega}\dot Q_i\,\mathrm dt
      &=-[\lambda_{Q_i}^{\omega}Q_i]
        +\int\dot\lambda_{Q_i}^{\omega}Q_i\,\mathrm dt.
\end{align}
Collecting the coefficients of $I_i$ and $Q_i$ then yields
\begin{align}
    \widetilde{\mathcal G}_i
    ={}&\sum_{m=0}^{M}\int_{t_m}^{t_{m+1}}
      \Big(
        \dot\lambda_{I_i}^{\omega}
        -\alpha\lambda_{I_i}^{\omega}
        +\omega_i\lambda_{Q_i}^{\omega}
        -\lambda_{Q_i}
      \Big)I_i\,\mathrm dt
      \nonumber\\
    &+\sum_{m=0}^{M}\int_{t_m}^{t_{m+1}}
      \Big(
        \dot\lambda_{Q_i}^{\omega}
        -\omega_i\lambda_{I_i}^{\omega}
        -\alpha\lambda_{Q_i}^{\omega}
        +\lambda_{I_i}
      \Big)Q_i\,\mathrm dt
      \nonumber\\
    &-\sum_{m=0}^{M}
      \left[
        \lambda_{I_i}^{\omega}I_i
        +\lambda_{Q_i}^{\omega}Q_i
      \right]_{t_m^+}^{t_{m+1}^-}.
    \label{eq:tes-omega-second-collected}
\end{align}
Choosing the two continuous-time coefficients to vanish and converting from
$\mathrm d/\mathrm dt$ to the reverse-time derivative gives precisely
Equations~\eqref{eq:tes-omega-helper-i} and
\eqref{eq:tes-omega-helper-q}.

\paragraph{Step 14: telescope the second-adjoint boundary terms.}
Let
\begin{equation}
    p_i(t)=\lambda_{I_i}^{\omega}(t)I_i(t)
           +\lambda_{Q_i}^{\omega}(t)Q_i(t).
\end{equation}
After the two interval coefficients have been cancelled,
Equation~\eqref{eq:tes-omega-second-collected} telescopes to
\begin{equation}
    \mathcal G_i
      =p_i(0^+)-p_i(T^-)
       +\sum_e\left[p_i(t_e^+)-p_i(t_e^-)\right].
    \label{eq:tes-omega-second-telescope}
\end{equation}
The oscillator transition has identity state Jacobian, so the secondary
adjoints are chosen continuous as in
Equation~\eqref{eq:tes-omega-helper-continuity}. Hence each event contributes
\begin{equation}
    p_i(t_e^+)-p_i(t_e^-)
      =\lambda_{I_i}^{\omega}(t_e)\Delta I_i
       +\lambda_{Q_i}^{\omega}(t_e)\Delta Q_i.
    \label{eq:tes-omega-second-event-term}
\end{equation}
The terminal helper values vanish, and the TeS oscillator starts from
$I_i(0)=Q_i(0)=0$, so both endpoint terms in
Equation~\eqref{eq:tes-omega-second-telescope} vanish. A threshold reset has
$\Delta I_i=\Delta Q_i=0$ and contributes nothing. At an arrival from
neuron $j$, however,
\begin{equation}
    \Delta I_i=0,
    \qquad
    \Delta Q_i=-w_{ij}.
\end{equation}
Therefore
\begin{equation}
    \frac{\mathrm d\mathcal L}{\mathrm d\omega_i}
      =-\sum_{(j,k)\to i}
        w_{ij}\lambda_{Q_i}^{\omega}(t^s_{jk}),
    \label{eq:tes-omega-gradient-derived}
\end{equation}
which is Equation~\eqref{eq:tes-omega-gradient}. The second augmentation is a
second adjoint stage, not a second-order differential equation: it replaces
the continuous state integral by one event-local sample at each arrival.

\paragraph{Why the impulse is applied to $Q$.}
The acceleration required in Equation~\eqref{eq:tes-slope-variation} can be
evaluated without numerical differentiation:
\begin{equation}
    a_{ik}
      =\frac{-\dot V^\theta_{ik}
              -\alpha I_i((t^\theta_{ik})^-)
              -\omega_iQ_i((t^\theta_{ik})^-)}{\tau_V}.
    \label{eq:tes-acceleration}
\end{equation}
Under the $Q$-injection in Equation~\eqref{eq:tes-arrival-jump}, $V_i$, $I_i$,
and $\dot V_i$ remain continuous at an arrival. The acceleration has only a
finite jump,
\begin{equation}
    \ddot V_i^+-\ddot V_i^-
      =\frac{\omega_iw_{ij}}{\tau_V},
    \label{eq:tes-acceleration-jump}
\end{equation}
so its one-sided values remain finite even when an input arrival coincides
with a threshold event. In contrast, direct $I$-injection would give
$\dot V_i^+-\dot V_i^-=w_{ij}/\tau_V$; consequently, $\ddot V_i$ would contain
the distributional term
$(w_{ij}/\tau_V)\delta(t-t^s_{jk})$ at the arrival. This is why TeS uses
quadrature injection when the delay depends on the threshold-crossing slope.

\paragraph{Regularity and scope.}
The resulting gradients are exact
where threshold crossings are transverse, the delay map is differentiable,
and the spike identities remain unchanged in a neighborhood of the evaluated
weight. As in
EventProp, the base factor $1/\dot V^\theta_{ik}$ remains singular at a grazing
crossing; making $\mathfrak d'_{ik}$ small suppresses only the additional
delay-dependent contribution. The frequency derivation additionally assumes
zero initial oscillator state and that the synaptic jump amplitude, reset map,
and delay law have no explicit dependence on $\omega_i$. An explicitly
frequency-dependent jump normalization would contribute an additional direct
event term.

The numerical gradient landscapes reported for this checkpoint compare the
implemented adjoint with centered finite differences only on locally stable
event topologies.  Agreement is a consistency and debugging result; it does
not by itself prove the analytical derivation above.
